\documentclass[11pt]{article}
\usepackage[a4paper,margin=1in]{geometry}
\usepackage{fontspec}
\usepackage[boldfont=false]{xeCJK}
\setCJKmainfont{FandolSong-Regular.otf}
\usepackage{amsmath}
\usepackage{amssymb}
\usepackage{array}
\usepackage{booktabs}
\usepackage{graphicx}
\usepackage{adjustbox}
\usepackage{enumitem}
\setlist[itemize]{topsep=2pt,partopsep=0pt,parsep=0pt,itemsep=2pt,leftmargin=1.4em}
\setlist[enumerate]{topsep=2pt,partopsep=0pt,parsep=0pt,itemsep=2pt,leftmargin=1.6em}
\usepackage{parskip}
\usepackage{fvextra}
\fvset{breaklines=true,breakanywhere=true,fontsize=\small}
\setlength{\parindent}{0pt}

\begin{document}

\begin{center}
  {\LARGE\bfseries Automated and emerging technologies}\\[0.35em]
  {\large Igcse Computer Science}
\end{center}
\vspace{0.6em}

\section*{Automated systems}

An \textbf{automated system} 自动化系统 is a mix of software and hardware that senses and responds to data in its \textbf{environment} 环境, with no need for \textbf{human intervention} 人工干预. In other words, it works on its own.

Examples of automated systems:

\begin{itemize}
\item a central heating 中央供暖 system;

\item a chemical process 化学过程 in a factory;

\item a \textbf{greenhouse} 温室 (for growing plants);

\item a \textbf{car park barrier} 停车场道闸.

\end{itemize}

\subsection*{Sensors, microprocessors and actuators}

Three parts work together in an automated system.

\begin{center}
\adjustbox{max width=\linewidth}{%
\begin{tabular}{ll}
\toprule
\textbf{Part} & \textbf{Job} \\
\midrule
\textbf{sensor} 传感器 & measures a physical quantity (such as temperature or light) and sends the data \\
\textbf{microprocessor} 微处理器 & compares the data with stored values and makes decisions \\
\textbf{actuator} 执行器 & receives a signal and causes movement or action (such as opening a valve) \\
\bottomrule
\end{tabular}}
\end{center}

The cycle works like this:

\begin{enumerate}
\item The sensor sends data to the microprocessor.

\item The microprocessor compares this data with \textbf{stored values} 存储值 and makes a decision.

\item The microprocessor sends signals to the actuators to take action.

\end{enumerate}

For example, in a greenhouse: a temperature sensor reads the heat; the microprocessor compares it with the wanted value; if it is too hot, the microprocessor signals an actuator to open a window.

\subsection*{Advantages and disadvantages}

Automated systems are used in many situations, such as industry, transport, agriculture 农业 (farming), weather (gathering data), gaming and lighting.

\begin{center}
\adjustbox{max width=\linewidth}{%
\begin{tabular}{ll}
\toprule
\textbf{Advantages} & \textbf{Disadvantages} \\
\midrule
work all day and night, without rest & cost a lot to buy and set up \\
faster and more \textbf{consistent} 一致的 than people & can break down and need expert repair \\
can work in places unsafe for people & may replace people's jobs \\
fewer human mistakes & cannot easily react to a situation they were not built for \\
\bottomrule
\end{tabular}}
\end{center}

\section*{Robotics}

\textbf{Robotics} 机器人技术 is the branch of technology that deals with the design, building and operation of \textbf{robots} 机器人.

A robot has these characteristics:

\begin{itemize}
\item a \textbf{physical structure} 物理结构 (a mechanical 机械的 body, such as arms);

\item \textbf{electrical components} 电子元件 (motors, sensors and wiring);

\item \textbf{programmable instructions} 可编程指令 — it follows a program that can be changed.

\end{itemize}

\subsection*{What robots do}

Robots can perform many roles, for example:

\begin{itemize}
\item \textbf{factory equipment} 工厂设备 — building cars, welding, moving heavy parts;

\item \textbf{domestic appliances} 家用电器 — such as a robotic \textbf{vacuum cleaner} 吸尘器.

\end{itemize}

\begin{center}
\adjustbox{max width=\linewidth}{%
\begin{tabular}{ll}
\toprule
\textbf{Advantages} & \textbf{Disadvantages} \\
\midrule
do dull or dangerous jobs & expensive to buy and maintain \\
work quickly and accurately & can replace human workers \\
do not get tired or bored & a robot has a \textbf{lack of independent decision-making} 缺乏独立决策 — it only does what it is programmed to do \\
\bottomrule
\end{tabular}}
\end{center}

A key limit of a robot is that it cannot think for itself. It cannot make its own choices when something unexpected happens.

\section*{Artificial intelligence}

\textbf{Artificial intelligence} 人工智能 (AI) is the \textbf{simulation} 模拟 of human intelligence by computer systems. This means a computer doing tasks that normally need human thinking.

\subsection*{Characteristics of AI}

AI systems usually have these features:

\begin{itemize}
\item they \textbf{collect data} and the \textbf{rules} 规则 for using that data;

\item they have the ability to \textbf{reason} 推理 (work things out using the rules);

\item they have the ability to \textbf{draw conclusions} 得出结论, which may be \textbf{approximate} 近似的 (a best guess) or \textbf{definite} 确定的 (certain);

\item they have the ability to \textbf{learn} 学习 from data;

\item they have the ability to \textbf{adapt} 适应 — to change their behaviour as they get new data.

\end{itemize}

\subsection*{Examples of AI}

\begin{itemize}
\item \textbf{expert systems} 专家系统 — software that gives advice like a human expert (for example, helping a doctor with a diagnosis);

\item \textbf{natural language processing} 自然语言处理 — understanding human speech or text;

\item \textbf{self-driving cars} 自动驾驶汽车.

\end{itemize}

\subsection*{Narrow, general and strong AI}

\begin{center}
\adjustbox{max width=\linewidth}{%
\begin{tabular}{ll}
\toprule
\textbf{Type} & \textbf{What it means} \\
\midrule
\textbf{narrow AI} 弱人工智能 & can do only one task or a small set of tasks (for example, a chess program). All AI today is narrow AI. \\
\textbf{general AI} 通用人工智能 & could do any task a human can do, switching between many different tasks. \\
\textbf{strong AI} 强人工智能 & would think and be aware like a real human mind. This does not exist yet. \\
\bottomrule
\end{tabular}}
\end{center}


\end{document}
